\chapter{Introduction}
\label{chap:intro}

This chapter motivates the thesis: it describes the workload (fleets
of LLM agents on HPC clusters) and the observability gap it creates,
states the observation the design builds on, and lists the
contributions and the document structure.

\section{Context: Multi-Agent LLM Systems on Clusters}
\label{sec:intro-context}

Applications built on large language models (LLMs) have changed shape
quickly over the last few years. The first generation consisted of single
request--response calls: a prompt in, a completion out. The second generation
wrapped the model in a loop, the \emph{agent}, in which the model is
prompted with a task and a transcript, emits a call to an external tool,
has the tool's result appended to that transcript, and is prompted again,
until it emits a final answer instead of a call~\cite{yao2023react}.
Section~\ref{sec:bg-agent-loop} gives the mechanism precisely; the term
\emph{agent} is used throughout this thesis for that loop and the process
running it, and for nothing more. The current generation
goes one step further and decomposes a task across \emph{multiple}
cooperating agents, each with its own role, context, and
conversation, which delegate sub-tasks to one another through structured
tool-call protocols such as the Model Context Protocol
(MCP)~\cite{anthropic2024mcp}. Multi-agent frameworks make this pattern
routine~\cite{wu2023autogen}: a planner agent delegates sub-tasks to worker
agents, workers consult specialist agents, and the overall computation is a
conversation \emph{between} programs as much as within them.

At the same time, the deployment platform of these systems is changing.
When agents were single processes, a laptop or a single cloud instance was
enough. Multi-agent workloads (scientific-workflow assistants, large-scale
code analysis, data-processing fleets) instead run on
\emph{clusters}: tens to hundreds of nodes, allocated through batch schedulers
such as SLURM~\cite{yoo2003slurm}, with several agents per node. Recent
work places LLM-agent orchestration directly on leadership-class HPC
systems~\cite{pham2026multiagenthpc} and on the workflow infrastructure
that spans them~\cite{souza2025provenance}. The result
is a genuinely distributed system in the classical sense: many processes on
many machines, communicating over a network, with no shared clock and no
single point of observation.

This combination is operationally demanding in a way neither ancestor was.
LLM agents are non-deterministic, expensive (every turn has a token cost),
slow (seconds per model call), and failure-prone in unfamiliar ways.
They retry, emit calls naming tools or arguments that do not exist, hang on
remote calls, and
silently degrade; empirical studies of multi-agent LLM systems find such
failures to be systematic rather than incidental~\cite{cemri2025mast}, and
the retry loops in particular are an instance of the amplification
pathologies long known to destabilize distributed
systems~\cite{bronson2021metastable}.
Running \emph{fleets} of them across a cluster multiplies these problems by
the number of nodes and adds the classical failure modes of distributed
infrastructure: slow nodes, network partitions, and stragglers that drag
down the whole computation~\cite{gunawi2018failslow}. An operator
needs to answer questions that span both layers at once: \emph{What is
the fleet doing right now? Which agents are talking to which? Why is this
scenario slow, and which hop is the bottleneck? Which node is dragging the
fleet down? Has this failure happened before?}

This setting is the object of study of this thesis: a fleet of LLM
agents distributed across the nodes of a high-performance-computing (HPC)
cluster, and the observability infrastructure required to operate it.

\section{The Gap: Agent Semantics vs.\ Cluster Topology}
\label{sec:intro-gap}

Observability tooling for this setting splits into two mature camps whose
intersection is, to our knowledge, unoccupied. The first camp,
\emph{LLM and agent observability} (LangSmith~\cite{langsmith},
Langfuse~\cite{langfuse}, Arize Phoenix~\cite{phoenix},
AgentOps~\cite{agentops}, the LLM products of established APM
vendors~\cite{datadogllm}, and the OpenTelemetry generative-AI semantic
conventions now consolidating them~\cite{otelgenai}), captures rich agent
semantics (prompts, tokens, costs, tool calls, rendered as traces of nested
spans in the lineage of Dapper~\cite{sigelman2010dapper}) but treats
the physical machines the agents run on as, at best, opaque metadata: host
identity is not a dimension the tools can group, compare, or alert over.
The second camp, \emph{HPC and fleet monitoring}
(Prometheus~\cite{prometheus} with Grafana over batch-scheduled clusters, and
HPC-native collectors such as LDMS~\cite{agelastos2014ldms}), is built for
that physical layer (per-node metrics at scale, over hundreds of
nodes) but carries no application semantics at all: that the process
saturating a node is an LLM agent mid-conversation, retrying a failed
delegation, is invisible. Section~\ref{sec:bg-sota} surveys both camps in
detail.

The consequence fits in one sentence, which this thesis takes as its
motivating problem:

\begin{quote}
\emph{The tools we survey answer ``what did my agent do?'' with a trace
of one run on one machine. To our knowledge, none answers ``what is
my agent \textbf{fleet} doing across the cluster, and which node is
dragging it down?''}
\end{quote}

Answering that question requires coupling agent semantics to physical
cluster topology in a single system, the intersection both camps leave
open.

\section{The Log-Based Approach}
\label{sec:intro-insight}

The technical starting point is an observation about the shape
of the data. Everything worth capturing about a multi-agent system is an
\emph{event}: an LLM turn (prompt, response, tokens, latency, cost), an
agent-to-agent call (caller, callee, correlation, outcome), a
retained-context operation (an item added to or evicted from the material
an agent's next prompt is assembled from, \S\ref{sec:bg-agent-signals}), a
recovery or backtrack. These events carry timestamps and causal
relationships: a delegation \emph{happens before} the turns it triggers,
which happen before the reply edge. A multi-agent execution, in other words,
\emph{is} a causally-ordered event log in the sense of
Lamport~\cite{lamport1978time}.

If the data is natively a distributed event log, the natural place to store
it is a \emph{distributed shared log}~\cite{balakrishnan2012corfu,
kreps2013log}, the abstraction industrialized by Kafka~\cite{kreps2011kafka}
and since carried into production control planes~\cite{balakrishnan2020delos}.
This thesis uses the term \emph{substrate} for that storage
system: the log store underneath the observability system, whichever one it
is. Agents append locally on the node where they run; the log
provides ordering, durability, and archival; and every observability view is
a \emph{reader} over the same substrate. Ordering and durability become
properties of the storage layer rather than of the instrumentation, and
cross-agent causal reconstruction comes with the substrate rather than being
bolted on.

This yields the hypothesis the thesis tests: \emph{if the events of a
multi-agent LLM fleet are captured into a distributed shared log whose store
places an append endpoint on every compute node, then agent semantics and
cluster topology are coupled by construction rather than by integration, and
a reader over that one log can answer the operator's fleet-scale questions
without further instrumentation.} It is testable in two parts.
Section~\ref{sec:problem-substrate} makes the storage half checkable by
stating the features a store must provide, so that any candidate can be
scored against it; Chapter~\ref{chap:evaluation} measures the second half
against the requirements of \S\ref{sec:problem-requirements}.

Which store provides the log is a separate question, and this thesis keeps
it separate: \S\ref{sec:problem-substrate}
states the contract a substrate must honour (five required features,
S1--S5, and four convenient ones, S6--S9, each of whose absence is priced
in a named mechanism), and Chapter~\ref{chap:design} is written against
that contract. The implementation runs on
ChronoLog~\cite{kougkas2020chronolog}, a distributed, tiered shared-log
store that orders events by physical time and, crucially for this setting,
places a log-capture process (a \emph{keeper}) on every cluster node, so
that the mapping from an event to the machine that produced it is a fact
about where the append landed rather than a field the instrumentation must
be trusted to fill in. That property, S6 in the contract, is what no
general-purpose commit log offers and what makes the coupling of agent
semantics to cluster topology structural. Section~\ref{sec:design-chronolog}
reports the ChronoLog-specific consequences in one place, so that the design
and its instantiation can be read, and criticized, separately.

On this foundation the thesis designs, implements, and validates an
observability framework for LLM-agent fleets on HPC clusters. It
is a research prototype: complete enough to run real agent fleets on a
live cluster and to be measured end to end, but validated at the scales
reported in Chapter~\ref{chap:evaluation} rather than hardened for
production use.
Two capture paths feed the log,
one for LLM turns and one for inter-agent
calls, and a dashboard reconstructs every view (live agent and cluster
topology, failure diagnosis, fleet health, cross-host tracing) purely
from what was captured. Because a shared log of this class is
optimized for ingest (archive reads trail appends by tens of seconds of
drain latency, and on the deployed client a live replay can block a reader
for minutes), the
system adopts a two-path read design: the log is the
durable \emph{cold} path, and an in-process \emph{hot} store serves real-time
views from the same event stream at ingest time. That split is the price of
one feature the substrate contract marks optional, and
\S\ref{sec:design-adapters} says precisely what a store providing it would
let a reimplementation delete. We validate the system
end-to-end with real Claude agents on live multi-node ChronoLog deployments.

\section{Contributions}
\label{sec:intro-contributions}

This thesis makes five contributions:

\begin{enumerate}
\item \textbf{A substrate contract for log-backed observability.} We
  identify the minimal set of features a storage substrate must provide for
  an observability system of this shape to be implementable on it (S1--S5),
  together with a second tier (S6--S9) whose absence is survivable at a
  stated price, and we score representative shared logs against both
  (Chapter~\ref{chap:problem}). The contract is what makes the rest of the
  design portable rather than particular: it separates the architecture from
  the workarounds one deployment required, and it lets a prospective adopter
  predict the cost of a port from a table rather than discover it during one.
\item \textbf{A two-path capture architecture over a distributed shared log.}
  LLM turns flow through a node-local \emph{spool}, an append-only file on
  the agent's own machine, drained into the log by a separate
  \emph{capture worker} process; inter-agent calls flow through a per-node
  \emph{collector} daemon into that node's keeper. The design provides
  at-least-once capture and effectively exactly-once ingest, deduplicating
  on a per-session sequence number (a \emph{high-water mark}), with crash
  recovery reconstructed from the log itself (Chapter~\ref{chap:design}).
\item \textbf{A hot/cold read design around the log's ingest-optimized
  trade-off.} The log store's drain latency (tens of seconds) and a
  replay path that can block the reader for minutes make it
  unreadable in real time; the system pairs it with an in-process
  \emph{hot store} written at the moment of ingest, and a
  server-sent-events (SSE) \emph{live bus} that pushes each event on to the
  browser, so that real-time views and durable
  full-fidelity views are projections of one event stream
  (Chapter~\ref{chap:design}).
\item \textbf{Fleet-scale analytics computed on the log.} Three analyses turn
  the captured events into operator answers: a failure-diagnosis pipeline,
  the \emph{Doctor}, which detects
  failures, groups the textually identical ones into ranked incidents, and
  attaches zero-cost diagnoses with cross-run incident history; fleet-health
  \emph{rollups}, per-(host, minute) buckets aggregated at write time, that
  make the cluster overview $O(\mathrm{nodes}\times\mathrm{buckets})$,
  independent of event volume; and cross-host critical-path tracing that
  names the bottleneck hop by \emph{self-time}, its latency minus its
  slowest child's (Chapter~\ref{chap:analytics}).
\item \textbf{A measured evaluation on live clusters.} An automated
  end-to-end harness exercises the full stack with real Claude Agent SDK
  agents on live 4- and 8-node ChronoLog clusters, and we measure
  five dimensions: capture overhead, freshness, detection quality against
  planted ground truth, scalability in event volume and agent count, and
  a timed cross-layer diagnosis case study
  (Chapter~\ref{chap:evaluation}).
  Real agents provide the end-to-end ground truth; a synthetic generator
  driving the \emph{identical} ingest path supplies the controlled,
  reproducible load for the scale sweeps, where thousands of turns of real
  inference would be neither affordable nor repeatable
  (\S\ref{sec:eval-method}).
  Section~\ref{sec:eval-limitations} bounds each of those claims: the
  scales actually validated, the diagnoser's coverage, and the substrate
  behaviour the design works around rather than fixes.
\end{enumerate}

\section{Document Structure}
\label{sec:intro-structure}

Chapter~\ref{chap:background} develops the background (distributed-systems
foundations, distributed tracing, the shared-log abstraction, ChronoLog, and
LLM agents) and surveys the state of the art to establish the gap.
Chapter~\ref{chap:problem} distills the target scenario into explicit
requirements and states the substrate contract S1--S9. The system design
follows in Chapter~\ref{chap:design}: the
capture paths, the mapping of agent semantics onto the log, delivery
semantics, and the hot/cold read split, written against that contract, with
\S\ref{sec:design-chronolog} collecting what the ChronoLog instantiation
adds. Chapter~\ref{chap:analytics} builds
the fleet-scale analyses on top of it. Chapter~\ref{chap:evaluation}
measures overhead, freshness, scalability, and end-to-end correctness, and
closes with a fault-localization case study. Chapter~\ref{chap:conclusions}
concludes.
